\chapter{Bumpy metrics, multiplicity, and abundance}\label{sec:bumpy-multiplicity-abundance}
In this section $M^{n+1}$ is closed. The width $\omega_p$ always denotes the mod-two volume spectrum defined in Definition \ref{def:p-width}; it is not an arbitrary fixed homotopy-class width.

\begin{definition}[Bumpy metrics]\label{def:bumpy-metric}
A smooth Riemannian metric is bumpy if every closed immersed minimal hypersurface has trivial kernel for its normal Jacobi operator. This includes finite covers; checking only embedded hypersurfaces is a weaker condition. See \cite{White17}, Theorem 2.1, and \cite{Zhou20}, Introduction.
\end{definition}

\begin{theorem}[Generic bumpiness]\label{thm:bumpy-metrics}
For a smooth closed manifold $M$, a Baire-generic smooth Riemannian metric has no closed immersed minimal submanifold with a nontrivial Jacobi field. Here generic means that the property holds on a residual set in the $C^\infty$ topology. This is \cite{White17}, Theorem 2.1 with the symmetry group taken to be trivial and $k=\infty$. The proof remains external.
\end{theorem}

\begin{theorem}[Multiplicity-one realization of the volume spectrum]\label{thm:zhou-multiplicity-one}
Let $(M^{n+1},g)$ be a smooth closed connected orientable Riemannian manifold with $2\leq n\leq6$ and a bumpy metric. For every positive integer $p$, there is a finite disjoint collection of smooth, connected, closed, embedded, two-sided minimal hypersurfaces $\Sigma_1,\ldots,\Sigma_q$ realizing
\[
 \omega_p(M,g)=\sum_{i=1}^q\operatorname{Area}_g(\Sigma_i),\qquad
 \sum_{i=1}^q\operatorname{index}(\Sigma_i)\leq p.
\]
Each component occurs with multiplicity one. We state an orientable specialization of \cite{Zhou20}, Theorem A and Theorem 5.2, retaining its mod-two volume-spectrum setting. This is an existence statement about a realization of $\omega_p$, not a claim about every critical varifold in an arbitrary min-max class. The proof uses prescribed-mean-curvature approximation and remains external.
\end{theorem}

\begin{theorem}[Infinitely many minimal hypersurfaces]\label{thm:song-infinitely-many}
Every smooth closed connected Riemannian manifold $M^{n+1}$ with $2\leq n\leq6$ contains infinitely many geometrically distinct smooth closed embedded minimal hypersurfaces. Neither bumpiness nor positive Ricci curvature is required. Source: \cite{Song18}, Theorem 1, in the supplied June 22, 2018 version. The statement concerns distinct hypersurfaces, not repeated integer multiples of one surface. Its proof remains external.
\end{theorem}

\begin{theorem}[Equidistribution for generic metrics]\label{thm:generic-equidistribution}
Let $M^{n+1}$ be smooth, closed, and connected, with $2\leq n\leq6$. For a $C^\infty$-generic metric $g$, there exists a sequence of smooth connected closed embedded minimal hypersurfaces $\Sigma_j$ such that, for every $f\in C^\infty(M)$,
\[
 \lim_{q\to\infty}
 \frac{\sum_{j=1}^q\int_{\Sigma_j}f\,dA_g}
      {\sum_{j=1}^q\operatorname{Area}_g(\Sigma_j)}
 =\frac{1}{\operatorname{Vol}_g(M)}\int_M f\,dV_g.
\]
This uses the total area of the first $q$ surfaces as the denominator, rather than the unweighted average of their separately normalized area measures. Source: \cite{MNS19}, Main Theorem, equation (1), p.~422. The theorem provides a suitable sequence; it does not assert equidistribution for every choice of width realizers. The proof remains external.
\end{theorem}
